\begin{titlepage}
\centering
{\large Universidade Federal da Paraíba\par}
{\large Centro de Informática\par}
\vspace{0.4cm}
{\large Introdução ao Processamento Digital de Imagens --- Período 2026.2\par}
{\large Professor: Leonardo Vidal\par}

\vspace{3.2cm}
{\LARGE\bfseries Reconstrução automática de sinais de áudio cardíacos a partir de traçados em imagem\par}
\vspace{0.8cm}
{\large Trabalho prático\par}

\vspace{3.5cm}
{\large Aluno: Alex Tavares Cordeiro\par}

\vfill
{\large Novembro de 2026\par}
\end{titlepage}

\begin{abstract}
\noindent Este trabalho reconstrói, a partir de uma imagem, três sinais de áudio de sons cardíacos: um sem sopro (Normal), um com o sopro inocente de Still e um com o sopro patológico de comunicação interventricular (VSD). A imagem é a Figura~1 do artigo de Tomas et al.\ (2011), que mostra os três sinais como gráficos em função do tempo. O programa, escrito em Python, isola cada gráfico, remove a grade, os eixos, os textos e as anotações, extrai a curva, calibra o eixo do tempo e o da amplitude, e grava os arquivos \texttt{normal.wav}, \texttt{still.wav} e \texttt{vsd.wav} (5000 amostras por segundo, 16 bits, mono). A filtragem espacial usa uma função de correlação bidimensional escrita do zero, no domínio do espaço, com preenchimento por zeros. A remoção da grade usa uma medida de cor que não depende do fundo, e a extração usa uma regra por coluna que preserva a amplitude dos picos. A análise dos sinais reconstruídos reproduz sete de oito afirmações do artigo: o tom do sopro de Still fica em 120~Hz, e o sopro do VSD tem banda larga. A figura entrega 1333 colunas por segundo, o que limita o conteúdo recuperável a cerca de 667~Hz.

\medskip
\noindent\textbf{Palavras-chave:} processamento digital de imagens; correlação espacial; extração de curvas; sons cardíacos; sopro de Still; comunicação interventricular.
\end{abstract}
\clearpage
